\section{Performance models and the limits of measured evidence}
\label{architecture-performance}

\subsection{Count work at the boundary that is actually limiting}

The Roofline model relates attainable arithmetic throughput to both compute
capacity and data movement. With operational intensity $I=W/B$, where $W$ is
floating-point work and $B$ is bytes transferred at a specified memory level,
its elementary bound is
\[
 P\leq\min(P_{\max},\beta I),
 \qquad
 t\geq\max\left(\frac{W}{P_{\max}},\frac{B}{\beta}\right).
\]
Here $\beta$ is sustained bandwidth at that same boundary. A memory-bound loop
does not reach twice the performance merely because arithmetic width doubles.
Williams, Waterman and Patterson develop this reasoning in
\href{https://www2.eecs.berkeley.edu/Pubs/TechRpts/2008/EECS-2008-134.pdf}{the
primary Roofline report}. It is a useful bound and diagnostic model, not a
prediction that every Nimbus loop reaches a roof.

For Nimbus, the choice of $B$ matters. Counting bytes requested by source
loads is different from measuring DRAM traffic after cache reuse. Host-to-device
traffic is a different boundary again, and shader access to resident VRAM has
another bandwidth. A CPU loop, a PCIe transfer, a GPU kernel and browser
texture upload should not be placed under one unexplained bandwidth number.
Extra validation or output scans can dominate a complete operation even when
its inner stencil has a favorable arithmetic intensity.

The end-to-end consequence follows by separating an accelerated fraction $f$
of runtime from everything else. If that fraction becomes $s$ times faster,
the idealized whole-program speedup is
\[
 S=\frac{1}{(1-f)+f/s}.
\]
This assumes the remaining work is unchanged and introduces no extra overhead.
If a kernel is 5 percent of total time, even eliminating it entirely saves at
most 5 percent under those assumptions. An actual change may add packing,
allocation, synchronization or validation, so the measured result can be worse.
The implication is to identify the costly owner before choosing an optimization,
not to optimize the largest-looking loop by inspection.

\subsection{Why a local speedup can disappear on another machine}

The cross-architecture investigations have exposed three distinct questions:
how much arithmetic an owner performs, how much data its consumers require,
and when the hardware can execute that work. These questions remain useful
without treating an early measurement on one host as a published benchmark or
an acceptance threshold for a class of computers.

For CPU execution, increasing the worker count can stop helping once a shared
cache, memory interface or synchronization point becomes limiting. SMT workers
share physical resources, and thread placement can change which memory and
cache paths an operator uses. Nimbus therefore exposes explicit worker and
OpenMP placement controls and records the execution context. Its topology
policy is a starting point. It does not encode a particular workstation's
best observed count as a universal x86-64 setting.

Pressure-material reuse addresses a different problem: successive acoustic
operations can need unchanged material coefficients while pressure itself
changes. The current owners retain the immutable information under explicit
state preconditions. That can remove repeated work in the pressure owner, but
the complete step also includes transport, boundaries, microphysics, reports
and synchronization. An owner-level improvement is a hypothesis about overall
performance until the full workload is compared. Similarly, removing an upload
is not sufficient if it adds another synchronization or shifts work to a less
favorable point in the dependency chain.

On a device with separate memory, a staged acoustic command includes upload,
execution, download and host publication. A timer around a transfer side can
include waits for preceding work, driver scheduling and visibility operations;
it is not automatically an isolated measurement of PCIe bandwidth. The staged
implementation keeps those phases explicit so they can be attributed and
compared. It must preserve the same boundary and failure semantics as the
shared-memory path before a transfer reduction is considered equivalent.

Power state, memory clocks, link state and competing work can vary between
runs of an identical binary. A monitor sampled at a different cadence from the
kernel trace cannot establish which individual call experienced a particular
state. Nimbus's measurement workflow therefore keeps uncertainty visible and
requires repeated, matched runs. The public repository does not publish a
personal-device baseline here. Current code documents the implemented reuse,
precision and memory-policy choices; broad ARM64, x86-64, NUMA and discrete-GPU
performance conclusions still require the corresponding measurements.

\subsection{A useful optimization experiment preserves the scientific job}

The candidate and baseline must perform the same work: same configuration,
duration, diagnostic cadence, output selection, memory policy, worker count and
admission requirements. Reducing scientific reports or returning fewer fields
can be a product-mode change, but is not an equivalent implementation speedup.
Likewise, a native-required profile cannot obtain an acceptable timing by
silently substituting a host operation.

Use an isolated owner measurement to locate cost, then evaluate the complete
runtime path with warmup identified and repeated baseline/candidate runs.
Alternating order helps reveal drift; a single before/after pair cannot
separate a small change from varying load or operating state. Nested timers
must state which intervals overlap. An acoustic sub-timer included in the
dynamics timer is not an additional independent contribution to total time.

Correctness evidence accompanies the timing: accepted field bytes or declared
numerical comparisons, checkpoints, relevant conservation ledgers, failures,
threading and applicable backend paths. A source-only count of removed
allocations is a hypothesis until the active runtime and retained dataflow
confirm it. A portable test on one generic x86-64 build demonstrates that
artifact's behavior there; it does not prove broad performance optimality.

\section{Maintaining the ownership and provenance boundaries}

Several inexpensive mistakes can invalidate a technically successful change.
A new source writer can omit a residency transition; a reused scratch array
can leave halos unreadably stale; a serializer can change canonical output
while catalog pins remain old; or a cancellation path can release a reservation
while a noninterruptible native call still owns its memory. These failures are
best understood by asking who owns each resource, which generation its values
describe, what constitutes completion, and which outputs must change together.

Maintained tests mirror those boundaries rather than relying on one aggregate
smoke test. SIMD/alias and hardware-policy tests cover CPU representation and
selection. Memory-plan and residency tests cover allocation and authority.
Native backend tests cover actual dispatch, rejection and completion. Output
and checkpoint tests cover publication and continuation. Studio tests cover
preflight, decoding, geometry, missing data, resource lifetime and pixels.
An analytic rendering fixture can expose an axis swap even when a storm image
looks plausible; a hash mutation can expose a stale contract even when a UI
control still renders successfully.

Generated files require their generator's exact contract. Formatting a
generated shader string header or an exact-scope header can break its text
verification without changing the apparent equation. Formatting a source file
whose raw hash is sealed changes that source identity even if serialized output
is unchanged. Such work should be isolated, regenerated and revalidated rather
than mixed into an unrelated scientific or repository cleanup. Archived
evidence must retain the identities it actually tested.

The repository's status vocabulary remains Engineering incomplete, Ready,
Campaign validated, Externally reviewed and Fully qualified. Source inspection,
bounded tests and this technical manual do not collapse those stages. In
particular, complete bookkeeping of an energy change is not proof that the
physical source is correct; a reproducible image is not proof of a physically
validated storm; and a successful one-host benchmark is not a general hardware
qualification.

\section{Primary references and how they are used}
\label{architecture-references}

The references below describe methods and platform semantics. Nimbus's actual
mode choices, ABI, implementation limits and admission remain defined by the
linked maintained source and tests. A paper's results are not measurements of
the current Nimbus executable.

\begin{description}
\item[Compression method.]
Peter Lindstrom (2014), \emph{Fixed-Rate Compressed Floating-Point Arrays},
IEEE Transactions on Visualization and Computer Graphics, 20(12), 2674--2683.
\href{https://doi.org/10.1109/TVCG.2014.2346458}{DOI: 10.1109/TVCG.2014.2346458}.
\href{https://vis.cs.ucdavis.edu/vis2014papers/TVCG/papers/2674_20tvcg12-lindstrom-2346458.pdf}{Primary paper PDF}.
Used for the block-compression design; Nimbus's containers and temporal
predictor are separately specified in the output writer.

\item[Compression modes and limits.]
ZFP 1.0.1 official documentation,
\href{https://zfp.readthedocs.io/en/release1.0.1/modes.html}{Compression Modes}
and \href{https://zfp.readthedocs.io/en/release1.0.1/faq.html}{FAQ}.
Used to distinguish fixed rate, precision, accuracy and reversible operation,
including finite-representation and pointwise-error caveats.

\item[Performance model.]
Samuel Williams, Andrew Waterman and David Patterson,
\emph{Roofline: An Insightful Visual Performance Model for Floating-Point
Programs and Multicore Architectures}, University of California, Berkeley
Technical Report UCB/EECS-2008-134 (2008).
\href{https://www2.eecs.berkeley.edu/Pubs/TechRpts/2008/EECS-2008-134.pdf}{Primary report PDF}.
Used as an operational-intensity bound, not as a fitted Nimbus performance curve.

\item[CPU instruction sets and compilation.]
Arm, \href{https://developer.arm.com/documentation/101028/0009/Advanced-SIMD--Neon--intrinsics}{Arm C Language Extensions: Advanced SIMD intrinsics};
Intel, \href{https://www.intel.com/content/www/us/en/developer/articles/technical/intel-sdm.html}{Intel 64 and IA-32 Architectures Software Developer Manuals};
GCC, \href{https://gcc.gnu.org/onlinedocs/gcc/x86-Options.html}{x86 Options}
and \href{https://gcc.gnu.org/onlinedocs/gcc/x86-Built-in-Functions.html}{x86 Built-in Functions}.
These define architectural and compiler semantics; emitted instructions and
Nimbus's compiled helper branches determine what a particular build uses.

\item[Thread placement.]
OpenMP Architecture Review Board.\par
\href{https://www.openmp.org/spec-html/5.1/openmpse62.html}{OMP\_PLACES}
and \href{https://www.openmp.org/spec-html/5.0/openmpse52.html}{OMP\_PROC\_BIND}.
Used to distinguish worker count, places and affinity. Runtime support and
observed placement remain necessary when interpreting a measurement.

\item[Vulkan ownership and execution.]
Khronos Group, \href{https://docs.vulkan.org/guide/latest/memory_allocation.html}{Memory Allocation},
\href{https://docs.vulkan.org/guide/latest/synchronization_examples.html}{Synchronization Examples},
and \href{https://docs.vulkan.org/spec/latest/chapters/pipelines.html}{Vulkan Pipelines}.
Used for the distinction between allocation types, memory visibility,
completion and graphics versus compute work.

\item[Metal resources.]
Apple, \href{https://developer.apple.com/documentation/metal/choosing-a-resource-storage-mode-for-apple-gpus}{Choosing a resource storage mode for Apple GPUs}.
Used for shared/private resource semantics. The current Nimbus backend's
resource and completion choices are documented separately above.

\item[Rendering framework.]
Kitware, \href{https://kitware.github.io/vtk-js/docs/structures.html}{VTK.js data structures}.
The pinned framework and Nimbus's workload policy, shader adaptations and
permanent rendering tests define actual Studio support.
\end{description}
